\section{Introduction}

Recent advances in Vision-Language-Action (VLA) models have established them as a promising framework for robotic manipulation, where pretrained Vision-Language Models (VLMs) are adapted to predict robot actions from visual observations and language instructions~\citep{brohan2022rt, zitkovich2023rt, black2410pi0, kim2025openvla, kim2025fine}. These models benefit from rich semantic and visual priors learned from internet-scale data, enabling promising performance across diverse manipulation tasks and environments~\citep{zhong2025survey}. Yet manipulation is fundamentally a dynamic process: success depends not only on recognizing \textit{what to do}, but also on anticipating \textit{what will happen} while doing it. However current policies typically predict action from the current observation reactively without explicitly modeling the future dynamics induced by these actions, limiting their robustness under out-of-distribution visual and spatial shifts~\citep{li2025surveyvisionlanguageactionmodelsembodied}.


A natural response to this limitation is to supervise the policy with \textbf{predictive objectives} that force it to anticipate the future, not just imitate the next action. Early efforts pursue this goal in the 2D image space, training the policy to forecast future RGB frames~\citep{cen2025worldvla, wu2024unleashing} or depth~\citep{zhangdreamvla} alongside action prediction. While such 2D signals are easy to obtain, they remain tied to appearance-level cues, with much of the supervision coming from static background context or nuisance visual variation rather than action-relevant changes~\citep{fei2025liberoplus, chen2025robotwin}. Recent works such as GeoPredict~\citep{qian2025geopredict} and Pri4R~\citep{kim2026pri4rlearningworlddynamics} address this limitation by supervising the policy with future 3D point tracks (4D), but at substantial cost: they either depend on dense tracks from external spatial trackers~\citep{xiao2025spatialtracker}, inflating preprocessing overhead~\citep{kim2026pri4rlearningworlddynamics}, or task the VLM with predicting tracks through extra queries~\citep{qian2025geopredict}, disrupting pretrained vision-language representations which may harm generalization~\citep{zhang2026vlmvla, driess2025knowledgeinsulatingvisionlanguageactionmodels}.


These limitations suggest three design principles for practical 4D supervision in VLA policies.  First, the 4D signal should be compact and easy to obtain. In many tabletop manipulation settings, much of the scene is static, while the most reliable and densely available motion signal comes from the robot embodiment itself. \emph{Future robot keypoint tracks}, the 3D trajectories of points anchored to the robot body, therefore provide a compact embodiment-centric 4D signal that can be computed from proprioceptive states via forward kinematics without external trackers. Second, this supervision should be injected without disrupting the pretrained VLM backbone or base policy, motivating lightweight auxiliary pathways with gradient isolation. Third, the 4D signal should be introduced at training-time only, leaving the policy's input-output interface unchanged at inference.

In this paper, we propose \textbf{ELAN4D}, a training framework that improves VLA policies with embodiment-centric 4D supervision via plug-and-play adaptation. As in Figure~\ref{fig:teaser}, we represent the robot body with a sparse set of 3D keypoints placed on the arm joints and end-effector, and derive their future displacement tracks from proprioceptive trajectories. These tracks provide metric, temporally dense 4D supervision without requiring external point tracking or dense scene reconstruction. To incorporate this supervision without disrupting the pretrained VLM or base policy, ELAN4D uses a plug-and-play ControlNet-style~\citep{zhang2023controlnet} auxiliary branch with a lightweight track decoder for predicting future robot keypoint displacements. This design preserves the backbone while allowing the learned 4D-aware features to support action prediction through a residual pathway. At inference time, the track decoder is discarded, preserving the base policy's input-output interface and requiring no additional queries or future predictions.

Through extensive experiments on LIBERO~\citep{liu2023libero}, LIBERO-Plus~\citep{fei2025liberoplus}, RoboTwin2.0~\citep{chen2025robotwin} and real-world manipulation tasks, we show that ELAN4D consistently improves base VLA across single-arm and bimanual tasks and performs favorably against SOTA VLA methods. Its improvements are particularly strong in out-of-distribution settings, including viewpoint, background, and layout shifts, demonstrating the effectiveness of embodiment-centric 4D supervision for policy learning.


In summary, our contributions are threefold. \textbf{First}, we introduce \textbf{ELAN4D}, a VLA framework for learning 4D-aware policies from future robot keypoint tracks. \textbf{Second}, we use these tracks as compact embodiment-centric 4D supervision and inject them through a ControlNet-style branch that preserves policies' inference interface. \textbf{Third}, we demonstrate that ELAN4D delivers consistent improvements over both simulation and real-world manipulation tasks, particularly in out-of-distribution settings requiring spatial generalization and visual robustness.

%In summary, our main contributions are threefold. First, we introduce \textbf{ELAN4D}, a VLA framework that equips VLA policies with 4D-aware representations through future robot keypoint track supervision. Second, we introduce robot keypoint tracks as compact, embodiment-centric supervision, and incorporate them through a ControlNet-style branch that preserves the policy interface at inference time. Third, we demonstrate that ELAN4D delivers consistent and substantial improvements over both simulation and real-world manipulation tasks, particularly in out-of-distribution settings requiring spatial generalization.
